% © 2026 Ing. David Jaroš — CC BY-NC-ND 4.0
%
% This work is licensed under a Creative Commons Attribution-NonCommercial-NoDerivatives
% 4.0 International License (CC BY-NC-ND 4.0).
%
% License History: Earlier drafts (up to v0.3) were released under CC BY 4.0.
% From v0.4 onward, all material is released under CC BY-NC-ND 4.0 to protect
% the integrity of the theoretical work during ongoing academic development.
%
% See LICENSE.md for full license text.
%
% Gap ID: Gap C1 (upgrade)
% Purpose: Provide geometric origin of the no-W_R rule from the ψ-circle
%          orientation and the matter sector assignment n>0.
%          Upgrades step3_gap_C1_resolution.tex from "model axiom" to
%          "geometric decoupling result" for SU(2)_R.
% Status: [MC] Motivated Candidate — decoupling derived; full algebraic exclusion open.
% Date: 2026-06-10

\ifdefined\INCLUDEMODE
\else
\documentclass[12pt,a4paper]{article}
\usepackage[a4paper,margin=2.5cm]{geometry}
\usepackage{amsmath,amssymb,amsthm}
\usepackage{mathtools}
\usepackage{hyperref}
\usepackage{xcolor}
\usepackage{mdframed}
\usepackage{booktabs}

\hypersetup{
  colorlinks=true,
  linkcolor=blue!60!black,
  citecolor=green!50!black,
  urlcolor=blue!60!black
}

\newtheorem{theorem}{Theorem}[section]
\newtheorem{lemma}[theorem]{Lemma}
\newtheorem{proposition}[theorem]{Proposition}
\newtheorem{corollary}[theorem]{Corollary}
\theoremstyle{definition}
\newtheorem{definition}[theorem]{Definition}
\theoremstyle{remark}
\newtheorem{remark}[theorem]{Remark}

\newmdenv[backgroundcolor=yellow!20, linecolor=orange, linewidth=1.5pt]{gapbox}
\newmdenv[backgroundcolor=green!15, linecolor=green!60!black, linewidth=1.5pt]{resultbox}
\newmdenv[backgroundcolor=red!10, linecolor=red!60, linewidth=1.5pt]{deadendbox}
\newmdenv[backgroundcolor=blue!8, linecolor=blue!50, linewidth=1.5pt]{insightbox}

\title{\textbf{Chirality Derivation, Step 4: Geometric Origin of the No-$W_R$ Rule}\\
       \large $SU(2)_R$ Decouples from $n>0$ Matter via $\psi$-Circle Parity}
\author{Ing. David Jaroš}
\date{June 2026}

\begin{document}
\maketitle

\begin{abstract}
We provide a geometric derivation of the no-$W_R$ rule, partially closing Gap~C1
from \texttt{step3\_gap\_C1\_resolution.tex}.
The argument proceeds in two steps.
First, the $\psi$-circle parity $P_\psi:\psi\to-\psi$ (Step~1) assigns
\emph{definite} parity to every weak-interaction current:
the observed left-chiral currents $J^{\mu,a}_L$ are $P_\psi$-\emph{odd},
whereas any hypothetical $SU(2)_R$ current $J^{\mu,a}_R$ built from
$n>0$ matter is $P_\psi$-\emph{even}.
Second, by the $P_\psi$-parity conservation of the UBT action (Step~1),
a $P_\psi$-even coupling $W_{R,\mu}\,J^{\mu,a}_R$ cannot be generated by
the same mechanism that generates the $P_\psi$-odd $W_\mu\,J^{\mu,a}_L$;
the two couplings live in disjoint parity sectors.
This upgrades Gap~C1 from a ``model axiom'' to a
\textbf{[MC] geometric decoupling result}:
$SU(2)_R$ does not couple to observed ($n>0$) matter at leading order
because the $\psi$-circle orientation assigns incompatible parity to the
corresponding currents.
A full algebraic exclusion of $SU(2)_R$ from the off-shell action
remains open (OP-S4).
\end{abstract}
\fi

%─────────────────────────────────────────────────────────────────────────────
\section{Prior Results and Chain Position}
\label{sec:prior}
%─────────────────────────────────────────────────────────────────────────────

\begin{center}
\begin{tabular}{lll}
  \toprule
  \textbf{Statement} & \textbf{Status} & \textbf{Reference} \\
  \midrule
  $P_\psi:\psi\to-\psi$ is a symmetry of $S[\Theta]$ & Proved [L1] &
    \texttt{step1\_psi\_parity.tex}, Lem.~2 \\
  $P_\psi$ acts as $\gamma^5$ on $\psi$-modes & Proved [L1] &
    \texttt{step1\_psi\_parity.tex}, Prop.~3 \\
  Matter lives in the $n>0$ sector & Proved [L1] &
    \texttt{step1\_psi\_parity.tex}, Lem.~4 \\
  $SU(2)_L$ from left action on $\mathcal{B}$ & Proved [L1] &
    \texttt{appendix\_E2\_SM\_geometry.tex}, §6 \\
  $W^\pm$ vertex is $P_\psi$-odd (conditional) & [MC] conditional &
    \texttt{step3\_gap\_C1\_resolution.tex} \\
  \textbf{$SU(2)_R$ decouples from $n>0$ sector} &
    \textbf{[MC] this file} & §\ref{sec:main} \\
  Full algebraic exclusion of $SU(2)_R$ & OPEN & OP-S4 \\
  \bottomrule
\end{tabular}
\end{center}

%─────────────────────────────────────────────────────────────────────────────
\section{Parity Structure of Gauge Currents}
\label{sec:parity}
%─────────────────────────────────────────────────────────────────────────────

\subsection{The $\psi$-parity operator}

From \texttt{step1\_psi\_parity.tex}:
\begin{equation}
  P_\psi:\psi\;\mapsto\;-\psi.
\end{equation}
On the mode expansion $\Theta = \sum_n \Theta_n(q,t)\,e^{in\psi/R_\psi}$:
\begin{equation}
  P_\psi\,\Theta_n \;=\; \Theta_n\,e^{-in\psi/R_\psi}.
\end{equation}
Thus $P_\psi$ maps winding number $n\to-n$.  For spinors assembled from the
$n>0$ sector (matter by Lemma~4 of Step~1), the $P_\psi$ action restricts to
the $\gamma^5$ eigenvalue structure on the even/odd modes.

\subsection{Parity of the left-chiral current}

The $SU(2)_L$ Noether current associated with the left-handed coupling is
(see \texttt{step3\_gap\_C1\_resolution.tex}, §3):
\begin{equation}
  J^{\mu,a}_L
  \;=\;
  \bar\Theta_L\,\gamma^\mu\,T^a\,\Theta_L,
  \quad
  T^a = \tfrac{i\sigma^a}{2}\cdot\mathbf{1}_R,
\end{equation}
where $\Theta_L$ denotes the $n>0$ (left-chiral) projection of $\Theta$.

\begin{lemma}[$J^{\mu,a}_L$ is $P_\psi$-odd]
\label{lem:JL_odd}
Under $P_\psi$:
\begin{equation}
  P_\psi\!\left(J^{\mu,a}_L\right) \;=\; -J^{\mu,a}_L.
\end{equation}
\end{lemma}

\begin{proof}
From Step~1, Proposition~3, $P_\psi$ acts as $\gamma^5$ on the $n>0$ spinor sector:
$P_\psi\Theta_L = \gamma^5\Theta_L$.
Under $P_\psi$:
\begin{equation}
  P_\psi(\bar\Theta_L)
  \;=\;
  \overline{P_\psi(\Theta_L)}
  \;=\;
  \overline{\gamma^5\Theta_L}
  \;=\;
  \bar\Theta_L\gamma^0(\gamma^5)^\dagger\gamma^0
  \;=\;
  -\bar\Theta_L\gamma^5,
\end{equation}
using $(\gamma^5)^\dagger = \gamma^5$ and $\gamma^0\gamma^5\gamma^0 = -\gamma^5$.
Therefore:
\begin{align}
  P_\psi(J^{\mu,a}_L)
  &= P_\psi(\bar\Theta_L)\,\gamma^\mu\,T^a\,P_\psi(\Theta_L) \notag\\
  &= (-\bar\Theta_L\gamma^5)\,\gamma^\mu\,T^a\,(\gamma^5\Theta_L) \notag\\
  &= -\bar\Theta_L\,\underbrace{\gamma^5\gamma^\mu\gamma^5}_{-\gamma^\mu}\,T^a\,\Theta_L \notag\\
  &= -J^{\mu,a}_L,
\end{align}
where $\gamma^5\gamma^\mu = -\gamma^\mu\gamma^5$ has been used.
\end{proof}

\subsection{Parity of a hypothetical right-handed current}

An $SU(2)_R$ current built from the same $n>0$ matter sector would take the form:
\begin{equation}
  J^{\mu,a}_R
  \;=\;
  \bar\Theta_{n>0}\,\gamma^\mu\,\tilde{T}^a\,\Theta_{n>0},
  \quad
  \tilde{T}^a = \mathbf{1}_L\cdot\tfrac{i\sigma^a}{2},
\end{equation}
where $\tilde{T}^a$ acts by \emph{right} multiplication:
$\tilde{T}^a M = M\cdot(i\sigma^a/2)$ for $M\in\mathcal{B}$.

\begin{lemma}[$J^{\mu,a}_R$ is $P_\psi$-even]
\label{lem:JR_even}
For the $SU(2)_R$ current built from $n>0$ sector matter,
\begin{equation}
  P_\psi\!\left(J^{\mu,a}_R\right) \;=\; +J^{\mu,a}_R.
\end{equation}
\end{lemma}

\begin{proof}
The right generator $\tilde{T}^a: M\mapsto M\cdot(i\sigma^a/2)$ commutes with
the left $\gamma^5$ action, since $\gamma^5$ acts on the \emph{left} matrix index
of $\mathcal{B}\cong\mathrm{Mat}(2,\mathbb{C})$ while $\tilde{T}^a$ acts on the \emph{right}:
\begin{equation}
  \gamma^5\!\left(\tilde{T}^a M\right)
  \;=\;
  \gamma^5 M \cdot \tfrac{i\sigma^a}{2}
  \;=\;
  \tilde{T}^a(\gamma^5 M).
\end{equation}
Therefore:
\begin{align}
  P_\psi(J^{\mu,a}_R)
  &= (-\bar\Theta\gamma^5)\,\gamma^\mu\,\tilde{T}^a\,(\gamma^5\Theta) \notag\\
  &= -\bar\Theta\,\gamma^5\gamma^\mu\,\tilde{T}^a\,\gamma^5\Theta \notag\\
  &= -\bar\Theta\,\gamma^5\gamma^\mu\gamma^5\,\tilde{T}^a\,\Theta \notag\\
  &= -\bar\Theta\,(-\gamma^\mu)\,\tilde{T}^a\,\Theta \notag\\
  &= +J^{\mu,a}_R.
\end{align}
In the third line we used $[\tilde{T}^a,\gamma^5]=0$ to move $\gamma^5$ past $\tilde{T}^a$.
\end{proof}

%─────────────────────────────────────────────────────────────────────────────
\section{Main Result: $SU(2)_R$ Geometric Decoupling}
\label{sec:main}
%─────────────────────────────────────────────────────────────────────────────

\begin{theorem}[$SU(2)_R$ decouples from $n>0$ sector — parity obstruction]
\label{thm:wr_decoupling}
Let the UBT action $S[\Theta]$ be invariant under $P_\psi$ (Step~1, Lemma~2).
Let matter be assigned to the $n>0$ sector (Step~1, Lemma~4).
The $P_\psi$-invariance of $S[\Theta]$ determines the parity of gauge fields
from their coupling to matter currents: the $W_\mu^a$ field must be $P_\psi$-odd
because it couples to the $P_\psi$-odd current $J^{\mu,a}_L$ (Lemma~\ref{lem:JL_odd}).
Under this constraint, any $SU(2)_R$ gauge field $W_{R,\mu}^a$ introduced with
the \emph{same} $P_\psi$-parity as $W_\mu^a$ (i.e., $P_\psi$-odd) cannot couple
consistently to $n>0$ sector matter without breaking $P_\psi$-invariance:
\begin{equation}
  P_\psi\!\left(W_{R,\mu}^a J^{\mu,a}_R\right)
  \;=\;
  \underbrace{(-W_{R,\mu}^a)}_{P_\psi\text{-odd}}
  \cdot\underbrace{(+J^{\mu,a}_R)}_{P_\psi\text{-even}}
  \;=\;
  -W_{R,\mu}^a J^{\mu,a}_R
  \;\neq\;
  W_{R,\mu}^a J^{\mu,a}_R.
\end{equation}
Therefore $W_{R,\mu}^a J^{\mu,a}_R$ cannot appear in a $P_\psi$-invariant action
when $W_{R,\mu}^a$ is constrained to be $P_\psi$-odd.
\end{theorem}

\begin{proof}
\textbf{Step 1: determine the parity of $W_\mu^a$ from the $SU(2)_L$ coupling.}
The left-handed coupling $W_\mu^a J^{\mu,a}_L$ is required to be $P_\psi$-invariant
(Step~1, Lemma~2).  Lemma~\ref{lem:JL_odd} gives $P_\psi(J^{\mu,a}_L) = -J^{\mu,a}_L$.
For $P_\psi(W_\mu^a J^{\mu,a}_L) = W_\mu^a J^{\mu,a}_L$:
\begin{equation}
  P_\psi(W_\mu^a)\cdot(-J^{\mu,a}_L) \;=\; W_\mu^a J^{\mu,a}_L,
\end{equation}
which requires $P_\psi(W_\mu^a) = -W_\mu^a$, i.e., $W_\mu^a$ is $P_\psi$-\emph{odd}.

\textbf{Step 2: test whether $SU(2)_R$ can couple via the same $P_\psi$-odd field.}
If $W_{R,\mu}^a$ is introduced with $P_\psi(W_{R,\mu}^a) = -W_{R,\mu}^a$
(the same parity as $W_\mu^a$), then for the coupling to be $P_\psi$-invariant
we need:
\begin{equation}
  P_\psi(W_{R,\mu}^a)\cdot P_\psi(J^{\mu,a}_R)
  \;=\;
  (-W_{R,\mu}^a)(+J^{\mu,a}_R)
  \;=\;
  -W_{R,\mu}^a J^{\mu,a}_R
  \;\neq\;
  W_{R,\mu}^a J^{\mu,a}_R,
\end{equation}
using Lemma~\ref{lem:JR_even}.
The coupling violates $P_\psi$-invariance.

\textbf{Step 3: alternative parity assignment.}
If instead $W_{R,\mu}^a$ is assigned $P_\psi$-even parity,
then $W_{R,\mu}^a J^{\mu,a}_R$ is $P_\psi$-even and would be
$P_\psi$-invariant in isolation.
However, this $P_\psi$-even $W_{R,\mu}^a$ would be a \emph{different} gauge field
from the $P_\psi$-odd $W_\mu^a$ — it cannot be identified with the Standard Model
$W^\pm$ field.  Such a field would couple exclusively to $P_\psi$-even (right-chiral)
currents, which belong to the $n<0$ sector (absent from observed matter by Lemma~4).
Hence it is observationally decoupled from all $n>0$ matter.
\end{proof}

\begin{insightbox}
  \textbf{Physical interpretation.}
  The $\psi$-circle orientation (AXIOM~B) breaks the
  $SU(2)_L\times SU(2)_R$ symmetry of $\mathcal{B}\cong\mathrm{Mat}(2,\mathbb{C})$
  by assigning matter to the $n>0$ winding sector.
  This sector carries $P_\psi$-\emph{odd} left-chiral currents.
  Any right-handed gauge coupling built from the same matter currents
  is $P_\psi$-\emph{even} and must be gauged by a $P_\psi$-\emph{even} field —
  but the $W^\pm$ gauge field is $P_\psi$-odd in the minimal theory.
  This is the geometric obstruction: $SU(2)_L$ and $SU(2)_R$ couplings
  live in \emph{different} $P_\psi$-parity sectors, and the minimal action
  selects only the $P_\psi$-odd sector.
\end{insightbox}

%─────────────────────────────────────────────────────────────────────────────
\section{What Remains Open}
\label{sec:open}
%─────────────────────────────────────────────────────────────────────────────

\begin{gapbox}
  \textbf{Open: full algebraic exclusion of $SU(2)_R$ (OP-S4).}

  Theorem~\ref{thm:wr_decoupling} shows that $SU(2)_R$ with the \emph{same}
  $P_\psi$-parity as $W^\pm$ cannot be added without breaking $P_\psi$
  invariance.
  However, it does not rule out:
  \begin{enumerate}
    \item A $P_\psi$-even $SU(2)_R$ gauge field that couples to $P_\psi$-even
          right-chiral currents in the $n<0$ sector (right-sector matter).
    \item A spontaneously broken $SU(2)_R$ that acquires a large mass and
          decouples from low-energy phenomenology.
    \item An $SU(2)_R$ that is algebraically consistent but whose mass
          is set to infinity by the $\psi$-circle radius $R_\psi\to 0$ limit.
  \end{enumerate}
  A complete no-$W_R$ theorem would require showing that \emph{no} consistent
  right-handed gauge coupling can be added to the off-shell UBT action,
  likely requiring the full algebraic structure of $\mathcal{B}$ and the
  compactness constraint on $\psi$.
\end{gapbox}

%─────────────────────────────────────────────────────────────────────────────
\section{Summary}
\label{sec:summary}
%─────────────────────────────────────────────────────────────────────────────

\begin{resultbox}
  \textbf{Gap C1 status after Step~4 (June 2026):}
  \begin{center}
  \begin{tabular}{lll}
    \toprule
    \textbf{Item} & \textbf{Before Step 4} & \textbf{After Step 4} \\
    \midrule
    No-$W_R$ rule & Model axiom & [MC] Geometric decoupling \\
    $W^\pm$ vertex is $P_\psi$-odd & Conditional & Conditional [MC] \\
    Full algebraic $SU(2)_R$ exclusion & Open & Open (OP-S4) \\
    \bottomrule
  \end{tabular}
  \end{center}

  The no-$W_R$ rule is now supported by a $\psi$-parity argument:
  $SU(2)_R$ couplings from $n>0$ matter live in the \emph{wrong parity sector}
  to be included in a $P_\psi$-invariant action alongside the $SU(2)_L$ coupling.
  This replaces the bare axiom by a structural motivation while leaving
  the full algebraic exclusion for future work.
\end{resultbox}

\ifdefined\INCLUDEMODE
\else
\end{document}
\fi
